\begin{summarybox}
	\textbf{\textcolor{CSummary}{一句话核心}}

	\textbf{从圆外一点引两条割线，两割线上从该点到两个交点的线段之积相等。}

	\medskip
	\textbf{\textcolor{CSummary}{使用条件}}
	\begin{itemize}[leftmargin=*, itemsep=0.5em, topsep=0.4em]
		\item \textbf{条件1（圆外一点）：} 点$P$在圆外。
		\item \textbf{条件2（两割线）：} 过点$P$的直线与圆有两个交点（即割线）。
	\end{itemize}

	\medskip
	\textbf{\textcolor{CSummary}{关键提醒}}
	\begin{itemize}[leftmargin=*, itemsep=0.5em, topsep=0.4em]
		\item \textbf{易错点（内外区分）：} 点$P$在圆内时改用相交弦定理，不可使用割线定理。
		\item \textbf{检查项（线段顺序）：} 乘积中的两条线段是从$P$到两个交点的距离，务必确认$PA$和$PB$（近点到远点），而非$AB$。
	\end{itemize}
\end{summarybox}
